\documentclass[10pt]{article}
\usepackage[letterpaper,margin=0.76in]{geometry}
\usepackage[T1]{fontenc}
\usepackage{lmodern}
\usepackage{microtype}
\usepackage{amsmath,amssymb,mathtools}
\usepackage{booktabs}
\usepackage{array}
\usepackage{xcolor}
\usepackage{titlesec}
\usepackage[font=small,labelfont=bf,skip=4pt]{caption}
\usepackage[numbers,sort&compress]{natbib}
\usepackage[most]{tcolorbox}
\usepackage{enumitem}
\usepackage{hyperref}

\definecolor{ink}{HTML}{1F2933}
\definecolor{muted}{HTML}{66737F}
\definecolor{accent}{HTML}{394B59}
\definecolor{boxbg}{HTML}{F2F5F7}
\hypersetup{colorlinks=true,linkcolor=accent,urlcolor=accent,citecolor=accent}
\color{ink}

\titleformat{\section}{\large\bfseries\color{accent}}{}{0pt}{}
\titlespacing*{\section}{0pt}{0.75em}{0.25em}
\setlength{\parindent}{0pt}
\setlength{\parskip}{0.40em}
\setlength{\bibsep}{1pt}
\renewcommand{\bibfont}{\small}

% --- Conceptual Persuasions callout -------------------------------------
\newtcolorbox{cpbox}{
  enhanced, breakable,
  colback=boxbg, colframe=boxbg,
  boxrule=0pt, arc=2.5pt, outer arc=2.5pt,
  left=10pt, right=10pt, top=7pt, bottom=8pt,
  borderline west={2.4pt}{0pt}{accent},
  before skip=0.95em, after skip=0.95em,
}
\newenvironment{persuasions}
  {\begin{cpbox}\small\setlength{\parskip}{0pt}%
   {\footnotesize\bfseries\color{accent}%
    \textls[110]{CONCEPTUAL PERSUASIONS}}\par\vspace{3pt}%
   \begin{itemize}[leftmargin=1.15em,itemsep=2.5pt,topsep=0pt,parsep=0pt,
                   label=\textcolor{accent}{\scriptsize$\blacktriangleright$}]}
  {\end{itemize}\end{cpbox}}

\title{\vspace{-2.3em}\textbf{Hubble Scale from Time-Induced Complexity}}
\author{J. R. Landers}
\date{}

\begin{document}
\maketitle
\vspace{-1.25em}

\begin{abstract}
The thinning of neighborhoods under added distinguishing structure suggests a
natural cosmological program.  As a three-dimensional world accumulates a
longer history, more relational information is required to specify its present
state; the thinning of compatible histories may then acquire a geometric
expression.  The simplest realization is unexpectedly suggestive: it predicts
$H_0=1/t_0=70.85\,\mathrm{km\,s^{-1}\,Mpc^{-1}}$ from the age of the Universe,
close to measured values.  This note develops the mathematical bridge from the
relational-shape theorem to that estimate and then uses cosmological data to
identify its next required ingredient.  Three spatial history channels imply
$a(t)\propto t$, while literal counting of three spatial dimensions plus time
implies $a(t)\propto t^{4/3}$.  These fixed laws are informative first
approximations; DESI DR2 baryon acoustic oscillation distances point beyond
them to an evolving rate of independent relational information.
A relativistic causal-past construction reinforces the same direction.  The
result is a concrete scale connection together with a sharply defined route
toward a covariant, dynamical model.
\end{abstract}

\section*{The relational-shape theorem and the proposed bridge}

The point of departure is the formalism of the companion paper, \emph{The
Relational Shape of Structural Complexity and Neighborhood Density}.  Let
\(
\mathcal O=\{O_1,\ldots,O_K\}
\)
be a population of $K\geq2$ objects.  At refinement depth $N$, object $O_i$ is
represented by $N$ metadata values,
\[
O_i^{(N)}=(o_{i1},\ldots,o_{iN}),
\]
where $i,j\in\{1,\ldots,K\}$ label objects and $n\in\{1,\ldots,N\}$ labels a
metadata coordinate.  The additive distance is
\begin{equation}
\delta_N(O_i,O_j)=\sum_{n=1}^{N}d_n(o_{in},o_{jn}),
\qquad d_n(x,y)\geq0,
\label{eq:additive-distance}
\end{equation}
where $d_n$ is the nonnegative dissimilarity contributed by coordinate $n$.
The mean separation of $O_i$ from the rest of the population is
\begin{equation}
D_N(i)=\frac{1}{K-1}\sum_{j\ne i}\delta_N(O_i,O_j).
\label{eq:mean-distance}
\end{equation}

For a tolerance radius $r\geq0$, define the neighborhood and its normalized
density by
\begin{align}
\mathcal N_r^{(N)}(O_i)
&=\{O_j\in\mathcal O:j\ne i,\ \delta_N(O_i,O_j)\leq r\},\\
\rho_N(i;r)
&=\frac{|\mathcal N_r^{(N)}(O_i)|}{K-1}
=\frac{1}{K-1}\sum_{j\ne i}
\mathbf 1\!\{\delta_N(O_i,O_j)\leq r\},
\label{eq:neighborhood-density}
\end{align}
where $|\cdot|$ denotes cardinality and $\mathbf 1\{\cdot\}$ is an indicator,
equal to one when its condition is true and zero otherwise.  Because the new
coordinate contribution is nonnegative,
\begin{equation}
\delta_{N+1}\geq\delta_N,
\qquad D_{N+1}(i)\geq D_N(i),
\qquad
\mathcal N_r^{(N+1)}(O_i)\subseteq\mathcal N_r^{(N)}(O_i),
\qquad \rho_{N+1}(i;r)\leq\rho_N(i;r).
\label{eq:thinning-theorem}
\end{equation}
The first inequality is shorthand for every pair $(O_i,O_j)$.  These
inequalities are the proved relational-shape result.  They require no
cosmology and no probabilistic assumption.  The associated discrete isolation
velocity,
\begin{equation}
v_i(N)\equiv D_{N+1}(i)-D_N(i)
=\frac{1}{K-1}\sum_{j\ne i}
d_{N+1}(o_{i,N+1},o_{j,N+1})\geq0,
\label{eq:isolation-velocity}
\end{equation}
measures the mean outward distance contributed by the newly appended
coordinate.

The cosmological analogy uses a particular probabilistic specialization.  Let
the coordinates be independent unbiased bits and let $d_n$ be Hamming
distance, $d_n(x,y)=\mathbf 1\{x\ne y\}$.  For the exact-match neighborhood
$r=0$, the ensemble-mean density is
\begin{equation}
\bar\rho_N\equiv\mathbb E[\rho_N(i;0)]
=\Pr[\delta_N(O_i,O_j)=0]=2^{-N},
\label{eq:binary-density}
\end{equation}
where $\mathbb E$ denotes expectation and $\Pr$ probability.  It is useful to
define the effective distinguishing information
\begin{equation}
C_N\equiv-\log_2\bar\rho_N.
\label{eq:complexity-information}
\end{equation}
Here $\log_2$ is the base-two logarithm.  Thus $C_N=N$ bits in this special
independent-binary model.  In general,
$C_N$, the coordinate count $N$, and the distance increment $v_i(N)$ need not
be equal: redundant coordinates can increase $N$ without adding a bit of
distinguishing information, while $v_i$ measures distance rather than the
logarithmic loss of neighbors.  For one unbiased binary coordinate,
$C_{N+1}-C_N=1$ bit whereas the ensemble-mean Hamming contribution is
$\mathbb E[v_i(N)]=1/2$.

Now comes the additional dictionary, not a consequence of the theorem.  Choose
a reference refinement $N_*$ and define inverse density as a dimensionless
relational-volume ratio,
\begin{equation}
\frac{V_{\rm rel}(N)}{V_{\rm rel}(N_*)}
\equiv\frac{\bar\rho_{N_*}}{\bar\rho_N}
=2^{C_N-C_{N_*}}.
\label{eq:relational-volume}
\end{equation}
Assigning that volume a three-dimensional linear scale gives
\begin{equation}
\frac{a_{\rm rel}(N)}{a_{\rm rel}(N_*)}
\equiv\left[\frac{V_{\rm rel}(N)}{V_{\rm rel}(N_*)}\right]^{1/3}
=2^{(C_N-C_{N_*})/3}.
\label{eq:relational-scale}
\end{equation}
If the discrete sequence $C_N$ admits a continuous interpolation $C(t)$ in
cosmic time $t$, an overdot denotes $d/dt$, $\ln$ is the natural logarithm,
and
\begin{equation}
H_{\rm rel}\equiv\frac{\dot a_{\rm rel}}{a_{\rm rel}}
=\frac{d\ln a_{\rm rel}}{dt}
=\frac{\ln2}{3}\dot C.
\label{eq:bridge}
\end{equation}
The earlier $N$-based formula is the special case $C=N$.  The complete logical
chain is therefore
\begin{equation}
\boxed{
\rho\ \longrightarrow\ C=-\log_2\rho
\ \longrightarrow\ V_{\rm rel}\propto2^C
\ \longrightarrow\ a_{\rm rel}\propto2^{C/3}
\ \longrightarrow\ H_{\rm rel}=\frac{\ln2}{3}\dot C .}
\label{eq:dictionary-chain}
\end{equation}

The mapping can be read directly against the relational formalism
(Table~\ref{tab:dictionary}):
\begin{center}
\captionof{table}{Proposed dictionary between the relational-shape quantities
and their cosmological readings.  The first five rows are internal to the
relational geometry; the last three are hypotheses added here.}
\label{tab:dictionary}
\small
\begin{tabular}{>{\raggedright\arraybackslash}p{0.29\linewidth}
                >{\raggedright\arraybackslash}p{0.62\linewidth}}
\toprule
Relational quantity & Proposed cosmological reading\\
\midrule
$O_i^{(N)}$ & A candidate state or history specified by $N$ accumulated constraints.\\
$\delta_N$ & Cumulative incompatibility between two candidate histories.\\
$\mathcal N_r^{(N)}(O_i)$ & Histories still compatible with $O_i$ within tolerance $r$.\\
$\rho_N(i;r)$ & Fraction of alternatives remaining locally compatible.\\
$v_i(N)$ & Outward distance added by the next distinguishing constraint.\\
$C=-\log_2\bar\rho$ & Effective independent distinguishing information, in bits.\\
$V_{\rm rel}\propto\bar\rho^{-1}$ & Hypothesized relational volume.\\
$a_{\rm rel}\propto V_{\rm rel}^{1/3}$ & Hypothesized three-dimensional scale factor.\\
\bottomrule
\end{tabular}
\end{center}
Only the first five rows belong to the relational geometry itself.  The last
three introduce the information measure and then identify its inverse density
with a volume and a scale.  Identifying $a_{\rm rel}$ with the physical
Friedmann--Lema\^itre--Robertson--Walker (FLRW) scale factor $a(t)$ is the first
specifically cosmological hypothesis.

\begin{persuasions}
\item Adding a new way to tell objects apart can only increase the distance
between them, never decrease it.  So neighborhoods shrink as detail
accumulates.  This much is proved.
\item If only one object in $2^C$ still matches you exactly, then $C$ is your
distinguishing information, measured in bits.
\item The added step: treat the number of remaining matches as a volume---the
fewer the matches, the larger the volume---and take its cube root to get a
length.
\item That length is then identified with the cosmic scale factor.  The
thinning is a theorem; this identification is an assumption.
\end{persuasions}

\section*{The accumulated-history idea}

A present state is not specified only by what exists now.  In a dynamical
world, it must also be compatible with an accumulated past.  As the temporal
baseline grows, two histories that once appeared equivalent may become
distinguishable.  The analogy with structural complexity is therefore direct:
history supplies additional relational constraints, thinning the set of states
that remain locally compatible with the present.

Operationally, fix a comoving domain---a spatial region carried with the mean
cosmic flow---and a comparison ensemble of candidate histories.  At cosmic
time $t$, $O_i^{(N(t))}$ denotes one history described by
the constraints accumulated so far, and coordinate $o_{in}$ records the value
of the $n$th independent relational constraint.  The radius $r$ specifies how
much cumulative incompatibility is still regarded as locally equivalent.  If
later times append rather than erase constraints, then for $t_2>t_1$ the
relational theorem gives
\begin{equation}
\mathcal N_r^{(N(t_2))}(O_i)
\subseteq\mathcal N_r^{(N(t_1))}(O_i),
\qquad
\bar\rho(t_2)\leq\bar\rho(t_1),
\qquad
C(t_2)\geq C(t_1).
\label{eq:history-filtration}
\end{equation}
Thus accumulated time is not inserted as a Euclidean distance coordinate.
Rather, time indexes a filtration of history neighborhoods, and $C(t)$ counts
the effective distinctions revealed by that filtration.  This is the direct
mathematical connection between accumulated history and structural
complexity.

The cleanest scale-free assumption is that each doubling of cosmic age adds a
fixed number $d_{\rm eff}$ of effective independent bits,
\begin{equation}
C(t)-C_*=d_{\rm eff}\log_2(t/t_*),
\label{eq:history}
\end{equation}
where $t_*>0$ is a reference cosmic time and $C_*=C(t_*)$.  In the independent
binary specialization one may write the same equation as
$N(t)-N_*=d_{\rm eff}\log_2(t/t_*)$, where $N_*=N(t_*)$.  Substitution into
Eqs.~\eqref{eq:relational-scale} and \eqref{eq:bridge}, followed by the
hypothesis $a=a_{\rm rel}$, yields
\begin{equation}
\frac{a(t)}{a_*}=\left(\frac{t}{t_*}\right)^p,\qquad
H(t)=\frac{p}{t},\qquad
p=\frac{d_{\rm eff}}{3}.
\label{eq:power}
\end{equation}
Here $a_*=a(t_*)$, $H\equiv\dot a/a$ is the physical Hubble rate, and $p$ is
the constant power-law index.  The hypothesis has become a definite
cosmology.  It is also now testable.

There are two closely related ways to interpret ``three dimensions plus
accumulated time,'' and it is important not to mix them.

In the first, history increases the information required to maintain a
\emph{three-dimensional} state.  One new distinction per spatial channel per
doubling gives $d_{\rm eff}=3$.  Then
\begin{equation}
C(t)-C_*=3\log_2(t/t_*),\qquad
a(t)\propto t,\qquad H(t)=\frac{1}{t}.
\label{eq:three}
\end{equation}
Time drives the accumulation, but is not itself counted as an interchangeable
fourth coordinate.

In the literal version, the three spatial directions and accumulated time are
all counted positively, so $d_{\rm eff}=4$.  This gives
\begin{equation}
C(t)-C_*=4\log_2(t/t_*),\qquad
a(t)\propto t^{4/3},\qquad H(t)=\frac{4}{3t}.
\label{eq:four}
\end{equation}
The distinction matters.  Equations~\eqref{eq:three} and \eqref{eq:four} make
different predictions and provide two clean benchmarks against which the
required evolution of $d_{\rm eff}$ can be measured.

\begin{persuasions}
\item A present state has to be consistent with everything that already
happened.  Each past event adds another constraint, so the list of constraints
grows with time.  Time supplies the coordinates; it is not one of them.
\item Two histories that looked identical early on become distinguishable
later, once enough record has built up to separate them.
\item The assumption: every doubling of the age of the Universe adds the same
number of new bits, $d_{\rm eff}$.  No epoch is treated as special.
\item That alone gives $a\propto t^{p}$ with $p=d_{\rm eff}/3$---a specific
expansion history, which data can check.
\item Three spatial channels means $d_{\rm eff}=3$ and no acceleration.
Counting time as a fourth means $d_{\rm eff}=4$ and slight acceleration.
\end{persuasions}

\section*{Why the present scale looks so promising}

Let $t_0=13.80\,\mathrm{Gyr}$ denote the measured present age of the Universe,
where $\mathrm{Gyr}$ means $10^9$ years, and let $H_0\equiv H(t_0)$ be the
present Hubble rate.  The superscripts $(3)$ and $(4)$ below label the
three- and four-channel hypotheses; they are not exponents.  The three-channel
law gives
\begin{align}
H_0^{(3)}&=\frac{1}{t_0}\notag\\
&=\frac{3.0856776\times10^{19}\,\mathrm{km/Mpc}}
{13.80\times3.15576\times10^{16}\,\mathrm{s}}\notag\\
&=70.85\,\mathrm{km\,s^{-1}\,Mpc^{-1}}.
\label{eq:h0-three}
\end{align}
Here $\mathrm{Mpc}$ denotes a megaparsec.  This lies $5.1\%$ above the Planck
base-$\Lambda$CDM value, where $\Lambda$CDM denotes the standard cosmology with
a cosmological constant $\Lambda$ and cold dark matter,
$67.4\pm0.5\,\mathrm{km\,s^{-1}\,Mpc^{-1}}$ \citep{Planck2018}, $3.9\%$
above the DESI DR2+cosmic microwave background (CMB) value used in the
accompanying analysis, $68.17\,\mathrm{km\,s^{-1}\,Mpc^{-1}}$
\citep{DESI2025}, and $3.0\%$ below the SH0ES distance-ladder value
$73.04\pm1.04\,\mathrm{km\,s^{-1}\,Mpc^{-1}}$ \citep{Riess2022}.
Because the CMB-inferred age and $H_0$ are strongly correlated and derived
within $\Lambda$CDM itself, Eq.~\eqref{eq:h0-three} is not an independent
prediction of $H_0$; it is a consistency of scale, not a measurement.

That is more than a dimensional order-of-magnitude match: it is a few-percent
match to the scale of the current expansion.  In information language, the
corresponding present rate is
\[
\dot C_0=\frac{3H_0}{\ln2}=0.3136\,\mathrm{bits/Gyr},
\]
where $\dot C_0=\dot C(t_0)$.  This is one effective distinction every $3.19$
billion years.  In the independent-binary specialization, $C=N$ and hence
$\dot C_0=\dot N_0$.

The literal four-channel law predicts a larger normalization:
\[
H_0^{(4)}=\frac{4}{3t_0}
=94.47\,\mathrm{km\,s^{-1}\,Mpc^{-1}}.
\]
Thus the especially close numerical connection belongs to the idea that
accumulated history increases the distinguishing complexity of a
three-dimensional state.  Treating time as a fourth Euclidean metadata
coordinate gives a distinct, readily testable benchmark.

The successful number also has a conventional cosmological explanation.  The
dimensionless product $H_0t_0$ is close to unity in the observed Universe:
approximately $0.95$ for Planck, $0.96$ for DESI DR2+CMB, and $1.03$ for SH0ES.
Any ansatz $H\sim1/t$ will therefore land near the current Hubble scale.  The
real test is whether it predicts the function $H(z)$ and the acceleration, not
whether it reproduces one inverse timescale today.

\begin{persuasions}
\item Put in the measured age, $13.80\,\mathrm{Gyr}$, and the model returns
$H_0=70.85$, within a few percent of every measurement.  Nothing was fitted.
\item The implied rate is very slow: about one new distinction per comoving
region every $3.2$ billion years.
\item But $H_0t_0$ is already close to $1$ in the real Universe, so any model
of the form $H\sim1/t$ would land here.  The agreement is expected rather than
informative.
\item The age itself was derived from $\Lambda$CDM, so comparing $1/t_0$ with
$\Lambda$CDM's own $H_0$ tests consistency, not prediction.
\item Matching one epoch says nothing about the rest of the expansion history.
That is what the next section tests.
\end{persuasions}

\section*{What the simplest construction teaches}

For a power-law scale factor $a(t)\propto t^p$, the deceleration parameter is
\begin{equation}
q=-\frac{a\ddot a}{\dot a^2}=\frac{1}{p}-1.
\end{equation}
Here $\dot a=d a/dt$ and $\ddot a=d^2a/dt^2$; negative $q$ denotes accelerated
expansion.
The three-channel model has $p=1$ and therefore $q=0$: it supplies the coasting
benchmark.  The literal $3+1$ model has $p=4/3$ and $q=-1/4$, so the same
formalism also naturally accesses accelerated expansion.  The observed
late-time value asks for a stronger and evolving effect.  For comparison, the
present flat-$\Lambda$CDM value for
$\Omega_m=0.3027$ is $q_0\simeq-0.546$, where $\Omega_m$ is the present matter
density as a fraction of the critical density and the subscript $0$ denotes
the present epoch.

The relational bridge itself suggests the generalization.  Define the
instantaneous effective number of new bits per doubling of cosmic age by
\begin{equation}
d_{\rm eff}(t)\equiv\frac{dC}{d\log_2 t}.
\label{eq:dynamic-deff}
\end{equation}
Equation~\eqref{eq:bridge} then gives the exact dictionary
\begin{equation}
d_{\rm eff}(t)=3H(t)t,
\qquad
H(t)=\frac{d_{\rm eff}(t)}{3t}.
\label{eq:dynamic-bridge}
\end{equation}
Consequently,
\begin{equation}
q(t)=-1-\frac{\dot H}{H^2}
=-1+\frac{3}{d_{\rm eff}(t)}
-\frac{3t\dot d_{\rm eff}(t)}{d_{\rm eff}(t)^2}.
\label{eq:dynamic-q}
\end{equation}
The constant-channel models set $\dot d_{\rm eff}=0$ and recover
$q=3/d_{\rm eff}-1=1/p-1$.  In the broader relational program, accelerated
expansion can depend both on the amount of new distinguishing information and
on how its rate changes.  The cosmological data can therefore be read as a
target curve for $d_{\rm eff}(t)$ rather than only as a verdict on a fixed
power law.

The expansion history supplies the stronger test.  Normalize $a(t_0)=a_0=1$
and define cosmological redshift by $1+z=a_0/a(t)$, so $z=0$ today and
$H_0=H(z=0)$.  Equation~\eqref{eq:power} predicts
\begin{equation}
H(z)=H_0(1+z)^{1/p}.
\label{eq:hz}
\end{equation}
Let $c$ be the speed of light and $r_d$ the comoving sound horizon at the
baryon-drag epoch.  In a spatially flat FLRW geometry, the radial Hubble
distance and transverse comoving distance are
\begin{equation}
D_H(z)\equiv\frac{c}{H(z)},\qquad
D_M(z)\equiv c\int_0^z\frac{dz'}{H(z')}.
\label{eq:bao-definitions}
\end{equation}
Writing the dimensionless nuisance scale as $X\equiv c/(H_0r_d)$ gives
\begin{align}
\frac{D_H(z)}{r_d}&=X(1+z)^{-1/p},\\
\frac{D_M(z)}{r_d}&=X\int_0^z(1+z')^{-1/p}\,dz'.
\end{align}
For $p=1$, the second expression is $X\ln(1+z)$.  For $p\ne1$ it is
\begin{equation}
\frac{D_M(z)}{r_d}
=X\frac{(1+z)^{1-1/p}-1}{1-1/p}.
\end{equation}
These are the complete distance--redshift predictions of the construction.

We compare them with the published 13-element DESI DR2 BAO vector and full
covariance \citep{DESI2025}.  The nuisance scale $X$ is fitted analytically for
every row, so the result isolates the redshift-dependent shape independently
of a model's chosen $H_0$ or sound horizon.
Writing the data vector as $\mathbf y$, a model prediction as $\mathbf m$, and
the published covariance matrix as $\boldsymbol\Sigma_{\rm BAO}$, the reported
statistic is
\begin{equation}
\chi^2=(\mathbf y-\mathbf m)^{\mathsf T}
\boldsymbol\Sigma_{\rm BAO}^{-1}(\mathbf y-\mathbf m),
\label{eq:chi-square}
\end{equation}
where $\mathsf T$ denotes transpose.  The degrees of freedom equal the number
of data values minus the number of fitted parameters.  The matter-like
$p=2/3$ row of Table~\ref{tab:bao} is included as a familiar power-law
benchmark.

\begin{center}
\captionof{table}{Fits to the 13-element DESI DR2 BAO vector with the nuisance
scale $X$ profiled analytically for every row.  Fixed power laws are compared
with the best single power and with flat $\Lambda$CDM.}
\label{tab:bao}
\begin{tabular}{lrr}
\toprule
History law & $\chi^2$ & degrees of freedom\\
\midrule
Matter-like, $p=2/3$ & 1420.18 & 12\\
Three spatial channels, $p=1$ & 187.09 & 12\\
Literal 3D plus time, $p=4/3$ & 1364.93 & 12\\
Best constant power, $p=0.9125$ & 96.97 & 11\\
Flat $\Lambda$CDM, $\Omega_m=0.3027$ & 10.63 & 12\\
\bottomrule
\end{tabular}
\end{center}

The comparison gives a precise design constraint.  The three-channel law
retains its striking present-day normalization, while the distance data show
that a constant information rate per logarithmic time cannot capture the full
redshift dependence.  Literal $3+1$ counting probes a different constant slope
and is still farther from the measured shape.  Even the best constant $p$
cannot follow all epochs.  Fitting the free power law only below $z=1$ gives
$p=1.1185$, but its six genuinely held-out predictions at $z\geq1$ have
$\chi^2=670.72$.  This is useful predictive information: the next construction
must promote $d_{\rm eff}$ to a dynamical quantity, or replace total accumulated
complexity by a conditional information rate.

\begin{persuasions}
\item A constant $d_{\rm eff}$ fixes the deceleration for all time: $p=1$ gives
none, $p=4/3$ gives a little.  The real Universe accelerates more than either.
\item Because $d_{\rm eff}(t)=3H(t)t$, any measured expansion rate converts
directly into a bit rate.  The data can be read in the model's own terms.
\item Acceleration can then come from two places: the size of $d_{\rm eff}$, or
its rate of change.  A fixed law only has the first.
\item No constant power fits the DESI distances.  Fitting $p$ below $z=1$ and
predicting the six points above it gives $\chi^2=670.72$---a clear failure of a
real prediction.
\item The failure is specific rather than fatal: it says $d_{\rm eff}$ must
change with time.
\end{persuasions}

\section*{Toward a relativistic information rate}

Time is not simply another spatial coordinate.  Proper time depends on the
worldline, causal order distinguishes past from future, and the spacetime
metric has Lorentzian rather than Euclidean signature.  Information-based
approaches that take spacetime seriously therefore work with local horizons or
causal order rather than appending a positive time distance
\citep{Jacobson1995,Surya2019}.

A natural relativistic extension is the four-volume of an observer's causal
past in flat FLRW spacetime,
\begin{equation}
V_4(\tau_o)=\frac{4\pi}{3}\int_0^{\tau_o}c\,d\tau'\,a^3(\tau')
\left[\int_{\tau'}^{\tau_o}\frac{c\,d\tau''}{a(\tau'')}\right]^3.
\label{eq:causal-volume}
\end{equation}
Here $\tau_o$ is the observer's proper time, $\tau'$ and $\tau''$ are integration
variables along the past history, $a(\tau)$ is the FLRW scale factor, and $c$
is retained so $V_4$ has dimensions of length to the fourth power.  The
observer has unit timelike four-velocity $u^\mu$, $\nabla_\mu$ is the spacetime
covariant derivative, and repeated Greek indices run over spacetime
coordinates.  The covariant expansion scalar is
$\theta\equiv\nabla_\mu u^\mu=3H$ for the comoving FLRW flow.

To test a constitutive connection, define either a raw causal-history variable
$\mathcal C_{\rm raw}=V_4$ or the dimensionless logarithmic variable
\begin{equation}
\mathcal C_{\log}(\tau_o)
\equiv\log_2\!\left[\frac{V_4(\tau_o)}{V_{4,*}}\right],
\label{eq:causal-complexity}
\end{equation}
where $V_{4,*}>0$ is an arbitrary reference four-volume.  If causal history
directly generated expansion with a constant constitutive coupling, one would
write
\begin{equation}
3H=\kappa\dot{\mathcal C},
\label{eq:constitutive-coupling}
\end{equation}
where the overdot here is $d/d\tau_o$ and $\kappa$ is the coupling between the
chosen causal-history variable and expansion.  Its units depend on whether
$\mathcal C$ is raw or logarithmic.  For the dimensionless bit count of the
direct relational dictionary, Eq.~\eqref{eq:bridge} would give
$\kappa=\ln2$.

Evaluated from $z=2.33$ to the present along the DESI-calibrated background,
the coupling required for raw $V_4$ changes by a factor of $321.8$.  Taking the
logarithm removes most of the drift, but the required coupling still changes by
a factor of $1.514$, with present value $0.5373$ rather than the $\ln2=0.6931$
suggested by the direct neighborhood dictionary.  This substantial
improvement from raw volume to logarithmic volume is encouraging: it shows
that information-like normalization is the right direction, while also
indicating that total causal volume alone is not yet the desired invariant.

The result points to a natural refinement.  Total causal volume counts shared
history again and again, whereas the relational-shape construction responds to
new distinctions.  A more faithful covariant quantity would therefore be the
conditional new information between neighboring causal slices after their
common past has been removed.  That extension preserves the core relational
idea while replacing accumulated amount by an independent information rate.

\begin{persuasions}
\item Time cannot be swapped for a spatial direction: elapsed time depends on
the path taken, and causal order separates past from future.  A covariant
version has to respect that.
\item The relativistic version of ``accumulated history'' is the four-volume of
the observer's causal past---everything that could have influenced them.
\item Using that volume directly fails: the coupling it requires changes by a
factor of $322$ since $z=2.33$.  Using its logarithm cuts the change to $1.5$.
\item The drift that remains has a clear cause.  The four-volume contains the
same early history at every step, so old information gets counted again and
again.
\item The fix is to count only what is new: the information one causal slice
adds beyond the one before it.
\end{persuasions}

\section*{The constructive next step}

The cosmological analogy makes quantitative contact immediately.  The exact
neighborhood law produces a transparent dictionary between distinction rate
and expansion rate, and the measured age of the Universe gives the correct
scale for $H_0$.  The distance--redshift curve then adds something valuable:
it shows that the relevant independent-information rate is dynamical rather
than a fixed channel count.

This leaves a concrete research program.  Accumulated history can be analogous
to increased structural complexity, with
$d_{\rm eff}(t)=dC/d\log_2t=3H(t)t$ providing the observational target.  The
next step is to define $C$ as an independent, covariant conditional-information
measure on neighboring causal slices and derive its dynamics.  Such a measure
would preserve the neighborhood-thinning mechanism while respecting the
special role of time and avoiding repeated counts of shared causal history.

The near equality $H_0\simeq1/t_0$ remains an elegant clue about why the idea
is quantitatively compelling: by Eq.~\eqref{eq:dynamic-bridge}, it says
$d_{\rm eff}(t_0)\simeq3$.  The simple model has therefore identified both the
natural present scale and the mathematical quantity whose evolution must be
explained.  That is a promising starting point for a relational account of
cosmic expansion.

\begin{persuasions}
\item What works is the chain from bit rate to expansion rate.  It reproduces
today's Hubble scale without being given it.
\item What fails is the assumption that $d_{\rm eff}$ is constant.  It has to
become a dynamical quantity.
\item Translated, $H_0\simeq1/t_0$ says only that $d_{\rm eff}\simeq3$ at
present.  It does not say three channels at every epoch.
\item The remaining work is well defined: build $C$ as the conditional
information between neighboring causal slices, derive how it evolves, and test
it against $d_{\rm eff}(t)=3H(t)t$, which the data already give.
\end{persuasions}

\bibliographystyle{unsrtnat}
\bibliography{../neighborhood_thinning_cosmology/complexity_cosmology}

\end{document}
